% Options for packages loaded elsewhere
\PassOptionsToPackage{unicode}{hyperref}
\PassOptionsToPackage{hyphens}{url}
\PassOptionsToPackage{dvipsnames,svgnames,x11names}{xcolor}
%
\documentclass[
  10pt,
  a4paper,
]{article}
\usepackage{amsmath,amssymb}
\usepackage{iftex}
\ifPDFTeX
  \usepackage[T1]{fontenc}
  \usepackage[utf8]{inputenc}
  \usepackage{textcomp} % provide euro and other symbols
\else % if luatex or xetex
  \usepackage{unicode-math} % this also loads fontspec
  \defaultfontfeatures{Scale=MatchLowercase}
  \defaultfontfeatures[\rmfamily]{Ligatures=TeX,Scale=1}
\fi
\usepackage{lmodern}
\ifPDFTeX\else
  % xetex/luatex font selection
\fi
% Use upquote if available, for straight quotes in verbatim environments
\IfFileExists{upquote.sty}{\usepackage{upquote}}{}
\IfFileExists{microtype.sty}{% use microtype if available
  \usepackage[]{microtype}
  \UseMicrotypeSet[protrusion]{basicmath} % disable protrusion for tt fonts
}{}
\makeatletter
\@ifundefined{KOMAClassName}{% if non-KOMA class
  \IfFileExists{parskip.sty}{%
    \usepackage{parskip}
  }{% else
    \setlength{\parindent}{0pt}
    \setlength{\parskip}{6pt plus 2pt minus 1pt}}
}{% if KOMA class
  \KOMAoptions{parskip=half}}
\makeatother
\usepackage{xcolor}
\usepackage[margin=22mm]{geometry}
\usepackage{listings}
\newcommand{\passthrough}[1]{#1}
\lstset{defaultdialect=[5.3]Lua}
\lstset{defaultdialect=[x86masm]Assembler}
\usepackage{longtable,booktabs,array}
\usepackage{calc} % for calculating minipage widths
% Correct order of tables after \paragraph or \subparagraph
\usepackage{etoolbox}
\makeatletter
\patchcmd\longtable{\par}{\if@noskipsec\mbox{}\fi\par}{}{}
\makeatother
% Allow footnotes in longtable head/foot
\IfFileExists{footnotehyper.sty}{\usepackage{footnotehyper}}{\usepackage{footnote}}
\makesavenoteenv{longtable}
\setlength{\emergencystretch}{3em} % prevent overfull lines
\providecommand{\tightlist}{%
  \setlength{\itemsep}{0pt}\setlength{\parskip}{0pt}}
\setcounter{secnumdepth}{-\maxdimen} % remove section numbering
\lstset{breaklines=true,basicstyle=\ttfamily\small}\setlength{\emergencystretch}{3em}
\ifLuaTeX
  \usepackage{selnolig}  % disable illegal ligatures
\fi
\IfFileExists{bookmark.sty}{\usepackage{bookmark}}{\usepackage{hyperref}}
\IfFileExists{xurl.sty}{\usepackage{xurl}}{} % add URL line breaks if available
\urlstyle{same}
\hypersetup{
  colorlinks=true,
  linkcolor={Maroon},
  filecolor={Maroon},
  citecolor={Blue},
  urlcolor={Blue},
  pdfcreator={LaTeX via pandoc}}

\author{}
\date{}

\begin{document}

\hypertarget{bpmlinearity-development-paired-conversion-gain-increment}{%
\section{BPM/linearity development: paired conversion gain
increment}\label{bpmlinearity-development-paired-conversion-gain-increment}}

Date: 3 November 2026. Status: exact numerical interface proposed for
owner review. Base: c73d9a9e3a631fbbb2346beaa3c150c110a60343. This
advances the accepted 2 November formula design; no new numerical source
is implemented by this review. Existing qualified bias/BPM execution
remains available.

\hypertarget{current-owner-decisions}{%
\subsection{Current owner decisions}\label{current-owner-decisions}}

\begin{itemize}
\tightlist
\item
  Continue BPM/linearity recipe development. QC investigation and
  visualization remain deferred.
\item
  The eventual linearity correction is applied in mefpp. Later
  validation reduces the SAME science frames with and without that
  correction and compares them in the photometric calibration notebook.
  That experiment is not this increment.
\item
  Use the selected Mac OmegaCAM sequence below. Nonconsecutive members
  are paired explicitly; chronological order and acquisition times
  remain unchanged.
\item
  Controlled lamps support testing small pair mismatches; do not infer
  perfectly constant illumination or introduce unrestricted per-exposure
  lamp parameters.
\item
  Keep the unsigned 8-bit reason BPM, independent coverage, reusable
  Eclipse numerical services, and the completed bias recipe
  orchestration template.
\end{itemize}

\hypertarget{source-reuse-review}{%
\subsection{Source reuse review}\label{source-reuse-review}}

Reviewed main's Eclipse sources and the existing BPM service:

\begin{longtable}[]{@{}
  >{\raggedright\arraybackslash}p{(\columnwidth - 4\tabcolsep) * \real{0.3333}}
  >{\raggedright\arraybackslash}p{(\columnwidth - 4\tabcolsep) * \real{0.3333}}
  >{\raggedright\arraybackslash}p{(\columnwidth - 4\tabcolsep) * \real{0.3333}}@{}}
\toprule\noalign{}
\begin{minipage}[b]{\linewidth}\raggedright
Existing resource
\end{minipage} & \begin{minipage}[b]{\linewidth}\raggedright
Actual role
\end{minipage} & \begin{minipage}[b]{\linewidth}\raggedright
Reuse decision
\end{minipage} \\
\midrule\noalign{}
\endhead
\bottomrule\noalign{}
\endlastfoot
src/iproc/image\_stats.c: image\_getstats\_masked\_checked & Finite,
BPM-selected ROI mean, median, sample standard deviation and optional
MAD & Reuse for four-frame means and counts \\
src/iproc/image\_stats.c: image\_pair\_stats\_checked & Pair difference,
intersected maps, checked regional statistics & Reuse for flat and
RAW-bias differences \\
src/iproc/detector.c: cube\_create\_gainmap\_robust / proportional &
Relative pixel response against measured frame level & Retain for
response BPM; not conversion gain in electrons/ADU \\
src/iproc/detector.c: detector\_linearity\_fit & Per-pixel time
polynomial in signal, no constant term; diagnostic is mean squared time
residual & Preserve legacy API; insufficient for the shared signed
shutter model \\
src/iproc/detector.c: cube\_correct\_detlin & In-place cubic using b/a
and c/a & Potential later correction primitive; explicit coefficient
conversion and validity checks required \\
src/math/fit\_curve.c and matrix.c & Polynomial/linear fits and
normal-equation least squares & Retain scaffolding; no joint constrained
covariance fit is supplied \\
\end{longtable}

The current BPM service already handles active ROIs and explicit sample
maps. File admission already carries pair identifiers. No replacement
FITS reader, statistics implementation, BPM response fitter or
acquisition-order sorting is needed.

\hypertarget{measurements-and-notation}{%
\subsection{Measurements and notation}\label{measurements-and-notation}}

The calculation uses \textbf{two flat-field exposures and two bias
exposures} from the same detector and readout configuration. A pair
means two exposures at the same integration time; they need not be
consecutive in the acquisition sequence.

An ROI (region of interest) is the rectangular detector area being
measured. The \textbf{same accepted pixels} must be used for every mean
and difference below. A pixel is accepted only when selected by the BPM
and finite in all four images. The recipe also excludes samples outside
the explicitly selected saturation and signal limits. There is no
automatic clipping.

\begin{longtable}[]{@{}
  >{\raggedright\arraybackslash}p{(\columnwidth - 4\tabcolsep) * \real{0.3333}}
  >{\raggedright\arraybackslash}p{(\columnwidth - 4\tabcolsep) * \real{0.3333}}
  >{\raggedright\arraybackslash}p{(\columnwidth - 4\tabcolsep) * \real{0.3333}}@{}}
\toprule\noalign{}
\begin{minipage}[b]{\linewidth}\raggedright
Symbol
\end{minipage} & \begin{minipage}[b]{\linewidth}\raggedright
Meaning
\end{minipage} & \begin{minipage}[b]{\linewidth}\raggedright
Unit
\end{minipage} \\
\midrule\noalign{}
\endhead
\bottomrule\noalign{}
\endlastfoot
\(p\) & Position of a pixel within the ROI & Pixel coordinate \\
\(G\) & Set of accepted pixel positions, common to all four images &
--- \\
\(N\) & Number of pixels in \(G\); at least two are needed for a
variance & Pixels \\
\(F_{1,p}\), \(F_{2,p}\) & Raw values at pixel \(p\) in the first and
second flat exposures & ADU \\
\(B_{1,p}\), \(B_{2,p}\) & Raw values at pixel \(p\) in the first and
second bias exposures & ADU \\
\(z_p\) & Any image or difference value at pixel \(p\) & ADU \\
\(\bar z\) & Arithmetic mean of \(z_p\) over \(G\) & ADU \\
\(v(z)\) & Sample variance of \(z_p\) over \(G\), with denominator
\(N-1\) & \(\mathrm{ADU}^{2}\) \\
\(\Delta F_p\) & Flat-pair difference, \(F_{1,p}-F_{2,p}\) & ADU \\
\(\Delta B_p\) & Bias-pair difference, \(B_{1,p}-B_{2,p}\) & ADU \\
\(g\) & Conversion gain: number of electrons corresponding to one ADU &
\(\mathrm{electrons}/\mathrm{ADU}\) \\
\(\mathrm{RON}_{\mathrm{ADU}}\) & Readout noise of one exposure,
expressed in ADU & ADU \\
\(\mathrm{MAD}(z)\) & Median absolute deviation from the median of
\(z_p\), over \(G\) & ADU \\
\end{longtable}

The subscript in \(\mathrm{RON}_{\mathrm{ADU}}\) is the complete unit
label ``ADU''; it is not a product of variables. Bars above image names,
such as \(\bar F_1\), mean averages over the common set \(G\).

\hypertarget{formulas-to-implement}{%
\subsection{Formulas to implement}\label{formulas-to-implement}}

\hypertarget{mean-and-sample-variance}{%
\subsubsection{Mean and sample
variance}\label{mean-and-sample-variance}}

For any image or difference \(z\):

\[
\bar z = \frac{1}{N}\sum_{p\in G} z_p ,
\qquad
v(z) = \frac{1}{N-1}\sum_{p\in G}(z_p-\bar z)^2 .
\]

The variance measures the scatter after removing the arithmetic mean.
All means and variances below use this same definition and pixel
population.

\hypertarget{readout-noise-from-the-raw-bias-pair}{%
\subsubsection{Readout noise from the raw bias
pair}\label{readout-noise-from-the-raw-bias-pair}}

Form the difference of the two raw biases:

\[
\Delta B_p = B_{1,p}-B_{2,p}.
\]

For independent exposures with equal readout-noise variance, the
difference contains twice the variance of one exposure. Therefore:

\[
\boxed{\mathrm{RON}_{\mathrm{ADU}}
       = \sqrt{\frac{v(\Delta B)}{2}} } .
\]

\textbf{No overscan correction is applied before this subtraction},
preserving the previously agreed raw-pair measurement without
scan-correction correlations.

The optional robust diagnostic is:

\[
\mathrm{MAD}(z)
 = \mathrm{median}_{p\in G}
   \left|z_p-\mathrm{median}_{q\in G}(z_q)\right| ,
\]

\[
\mathrm{RON}_{\mathrm{MAD}}
 = \frac{1.4826}{\sqrt{2}}\,\mathrm{MAD}(\Delta B).
\]

Here \(q\) is just a second pixel index used to express the inner
median. The factor 1.4826 puts MAD on the standard-deviation scale for a
Gaussian distribution. The MAD estimate is reported separately; it does
not replace ordinary variance in the conversion-gain calculation.

\hypertarget{conversion-gain-from-flat-and-bias-pairs}{%
\subsubsection{Conversion gain from flat and bias
pairs}\label{conversion-gain-from-flat-and-bias-pairs}}

Form the raw flat-pair difference:

\[
\Delta F_p=F_{1,p}-F_{2,p}.
\]

The conversion gain is:

\[
\boxed{
g =
\frac{\bar F_1+\bar F_2-\bar B_1-\bar B_2}
     {v(\Delta F)-v(\Delta B)}
}
\quad [\mathrm{electrons}/\mathrm{ADU}] .
\]

The numerator is the sum of the two flat signals after removing the
estimated bias levels. The denominator removes the readout-noise
contribution from the flat-difference variance. The relation assumes
photon shot noise, compatible flat illumination patterns and operation
within the accepted linear signal range. Both numerator and denominator
must be positive for a reported gain.

Use raw flat differences over the same accepted pixels, without scan
correction, for this first photon-transfer estimator. Bias removal in
the numerator uses the four image means. Changes in electronic structure
between exposures remain an admission concern; this routine does not
model or correct arbitrary drift.

\hypertarget{variance-versus-signal-which-value-represents-the-instrument}{%
\subsubsection{Variance versus signal: which value represents the
instrument?}\label{variance-versus-signal-which-value-represents-the-instrument}}

There is no single instrument variance that applies at every signal
level. Distinguish \textbf{readout variance at zero signal} from the
\textbf{total variance at a specified signal}, and distinguish both from
conversion gain.

For flat pair \(j\), define the bias-subtracted average signal and half
the difference variance:

\[
\mu_j =
\frac{\bar F_{j1}+\bar F_{j2}-\bar B_1-\bar B_2}{2},
\qquad
V_j=\frac{v(F_{j1}-F_{j2})}{2}.
\]

Here \(j\) labels an integration-time pair, and subscripts \(j1\) and
\(j2\) identify its two flat exposures. \(\mu_j\) is in ADU and \(V_j\)
is in squared ADU. The factor of two converts the difference variance to
the variance of ONE exposure, assuming independent exposures with
comparable noise statistics. For modest illumination mismatch it
represents the average of their variances.

Define the independently measured raw-bias readout variance:

\[
V_0=\frac{v(\Delta B)}{2}
   =\mathrm{RON}_{\mathrm{ADU}}^{\,2}.
\]

In the constant-gain, independent-pixel, photon-noise regime:

\[
\boxed{ V(\mu)=V_0+\frac{\mu}{g_0} } .
\]

The new symbol \(g_0\) denotes a reference conversion gain in
electrons/ADU, measured over an explicitly selected approximately linear
signal interval. \(V(\mu)\) denotes the single-exposure variance at
bias-subtracted signal \(\mu\); \(V_0\) is its zero-signal readout
component.

If plotting the FULL difference variance instead, as in the owner's
question:

\[
v(\Delta F)=2V_0+\frac{2\mu}{g_0}.
\]

Thus the slope of half-difference variance versus average signal is
\(1/g_0\); the slope of full-difference variance versus average signal
is \(2/g_0\). The ordinate is NOT the spatial variance of the average
flat image.

\textbf{Proposed reporting policy, for review:}

\begin{itemize}
\tightlist
\item
  Use \(V_0\) from the raw bias pair for readout variance. Zero signal
  is the reference; no illuminated flat point is selected to define
  readout noise.
\item
  Estimate \(g_0\) over a declared, verified approximately linear signal
  interval, above the region where read noise dominates and below
  significant curvature or saturation. Use a range of accepted pairs,
  not one arbitrarily chosen point. Record the interval, pair
  membership, fitting method and uncertainty; no numerical interval or
  tolerance is selected in this proposal.
\item
  Keep the individual pair estimates versus signal:
\end{itemize}

\[
  g_{\mathrm{pair},j}=\frac{\mu_j}{V_j-V_0}.
\]

In a curved relation this is a signal-dependent effective estimate. It
need not equal a gain inferred from the local tangent. A local-slope
diagnostic, if later implemented with a reviewed curve model, would be

\[
  g_{\mathrm{slope}}(\mu)
  =\left(\frac{dV}{d\mu}\right)^{-1}.
\]

Both quantities agree with \(g_0\) for a straight photon-transfer
relation; they generally differ for a curved one. Neither alone proves
an intrinsic electronics gain change.

\begin{itemize}
\item
  For the variance needed at a particular science signal, evaluate the
  validated variance-versus-signal model at that signal. Do not assign
  \(V_0\) or the variance at a calibration flat level to every science
  pixel.
\item
  If no sufficiently constrained approximately linear interval exists,
  report the reference gain as undetermined and retain the measured
  curve and pair estimates. Do not force a constant gain or extrapolate
  from saturation.
\end{itemize}

The intended first helper still returns each pair's measured quantities.
A fitted reference gain/curve belongs to subsequent sequence-level
analysis; this reporting proposal does not silently add a new curve
fitter to that helper.

Curvature can reflect nonlinear response, changing pixel
correlations/charge redistribution, illumination mismatch or other
signal-dependent effects. Consequently, a straight-line fit across the
whole curve cannot establish a universal gain, and a per-pixel variance
alone does not describe covariance. The later application of a linearity
correction must propagate the noise through that correction, as required
by the controlling design.

As a methodological reference,
\href{https://www.eso.org/observing/dfo/quality/VIMOS/qc/detlin_qc1.html}{ESO's
VIMOS gain documentation} describes paired estimates and photon-transfer
fitting, and records a later restriction to low-signal flats. That
instrument-specific threshold is not adopted for OmegaCAM.

\hypertarget{median-centering-and-controlled-lamps}{%
\subsubsection{Median centering and controlled
lamps}\label{median-centering-and-controlled-lamps}}

Let \(D_p=\Delta F_p\) and let \(m\) be its median over \(G\):

\[
m=\mathrm{median}_{p\in G}(D_p),
\qquad D'_p=D_p-m.
\]

Subtracting this one constant makes the median of the difference zero,
but:

\[
v(D')=v(D),
\qquad
\mathrm{MAD}(D')=\mathrm{MAD}(D).
\]

Thus your proposed centering changes the displayed zero point, not the
measured variance or MAD. No extra centering estimator or input mutation
is needed. Retain the original difference median and the two
bias-subtracted flat levels as pair-consistency diagnostics. A
multiplicative lamp change can leave a spatial response pattern after
median centering; no automatic rescaling is added. The original
uncentered levels remain available for the later linearity fit.

\hypertarget{proposed-eclipse-interface}{%
\subsection{Proposed Eclipse
interface}\label{proposed-eclipse-interface}}

This is the C entry point for the measurements defined above:

\begin{lstlisting}[language=C]
int detector_gain_pair_checked(
    const image_t *flat_first, const image_t *flat_second,
    const image_t *bias_first, const image_t *bias_second,
    const pixelmap *selection, const int *zone,
    int with_mad, detector_gain_pair_result *out);
\end{lstlisting}

\begin{itemize}
\tightlist
\item
  The four image arguments supply the flat pair and the raw bias pair.
\item
  \passthrough{\lstinline!selection!} is the optional common good-pixel
  map in Eclipse's existing convention: 1 means accepted, 0 means
  excluded.
\item
  \passthrough{\lstinline!zone!} supplies the ROI boundaries, inclusive
  and counted from 1.
\item
  \passthrough{\lstinline!with\_mad!} requests the additional robust
  diagnostic.
\item
  \passthrough{\lstinline!out!} receives the measurements, counts and
  availability/reason fields.
\end{itemize}

All four inputs have the same detector/readout, binning, dimensions and
raw-to-active mapping. The recipe uses the existing crop operation to
provide active-pixel images without scan subtraction. The selection is
an explicit binary Eclipse view, not the eight-bit reason array.

The result retains the four image-level statistics, both difference
statistics, the common valid/excluded counts, bias-subtracted pair
signals, pair-mismatch diagnostics, conversion gain and raw-pair RON.

\hypertarget{errors-ownership-and-callers}{%
\subsubsection{Errors, ownership and
callers}\label{errors-ownership-and-callers}}

Inputs/maps/ROI are borrowed and immutable; output is caller-owned and
has no owned image allocations. Temporary images/maps use existing
Eclipse allocation, cleanup and error conventions. Invalid
dimensions/maps/ROI, invalid options, allocation failure or nonfinite
arithmetic return an error without changing *out. Fewer than two common
samples or nonpositive gain numerator/denominator produce an explicit
unavailable measurement, never a nominal zero gain. RON can remain
available when gain is unavailable. Reject aliasing/invalid output
arguments as specified in the matching header documentation.

The numerical routine reports measurements, not scientific acceptance.
The recipe applies explicit illumination-mismatch/range thresholds and
records the decision. It calls the helper once per configured
amplifier/ROI/pair and records detector identity, sample population, raw
origin and estimator. Overlapping ROI estimates share pixels and are not
independent observations.

No new global model fitter, default mismatch tolerance, defect criterion
or uncertainty approximation is approved by this interface review.

\hypertarget{a-small-synthetic-example-for-checking-the-arithmetic}{%
\subsection{A small synthetic example for checking the
arithmetic}\label{a-small-synthetic-example-for-checking-the-arithmetic}}

``Four-pixel fixture'' meant \textbf{four artificial images, each
containing only four pixels}. It did not mean four real OCAM exposures
or a recommended measurement area. The four pixel positions are the same
in every artificial image. All sixteen values below are invented for a
simple software test.

\begin{longtable}[]{@{}
  >{\raggedright\arraybackslash}p{(\columnwidth - 8\tabcolsep) * \real{0.1579}}
  >{\raggedleft\arraybackslash}p{(\columnwidth - 8\tabcolsep) * \real{0.2105}}
  >{\raggedleft\arraybackslash}p{(\columnwidth - 8\tabcolsep) * \real{0.2105}}
  >{\raggedleft\arraybackslash}p{(\columnwidth - 8\tabcolsep) * \real{0.2105}}
  >{\raggedleft\arraybackslash}p{(\columnwidth - 8\tabcolsep) * \real{0.2105}}@{}}
\toprule\noalign{}
\begin{minipage}[b]{\linewidth}\raggedright
Pixel position \(p\)
\end{minipage} & \begin{minipage}[b]{\linewidth}\raggedleft
First flat \(F_{1,p}\)
\end{minipage} & \begin{minipage}[b]{\linewidth}\raggedleft
Second flat \(F_{2,p}\)
\end{minipage} & \begin{minipage}[b]{\linewidth}\raggedleft
First bias \(B_{1,p}\)
\end{minipage} & \begin{minipage}[b]{\linewidth}\raggedleft
Second bias \(B_{2,p}\)
\end{minipage} \\
\midrule\noalign{}
\endhead
\bottomrule\noalign{}
\endlastfoot
1 & 101 & 99 & 11 & 10 \\
2 & 99 & 101 & 9 & 10 \\
3 & 100 & 100 & 10 & 10 \\
4 & 100 & 100 & 10 & 10 \\
\end{longtable}

All entries are in ADU. All four positions are accepted, so \(N=4\). The
four image means are:

\[
\bar F_1=\bar F_2=100\ \mathrm{ADU},
\qquad
\bar B_1=\bar B_2=10\ \mathrm{ADU}.
\]

The pair differences, in pixel order, are:

\[
\Delta F=(2,-2,0,0)\ \mathrm{ADU},
\qquad
\Delta B=(1,-1,0,0)\ \mathrm{ADU}.
\]

Both difference means are zero. With the sample-variance denominator
\(N-1=3\):

\[
v(\Delta F)=\frac{2^2+(-2)^2+0^2+0^2}{3}
           =\frac{8}{3}\ \mathrm{ADU}^{2},
\]

\[
v(\Delta B)=\frac{1^2+(-1)^2+0^2+0^2}{3}
           =\frac{2}{3}\ \mathrm{ADU}^{2}.
\]

The expected software-test results are therefore:

\[
g=\frac{100+100-10-10}{8/3-2/3}
  =90\ \mathrm{electrons}/\mathrm{ADU},
\]

\[
\mathrm{RON}_{\mathrm{ADU}}
 =\sqrt{\frac{2/3}{2}}
 =\sqrt{\frac{1}{3}}\ \mathrm{ADU}
 \simeq0.57735\ \mathrm{ADU}.
\]

These deliberately simple numbers make errors in subtraction, the
\(N-1\) denominator or the factor of two easy to detect. \textbf{They
are not expected OCAM gain/noise values and do not validate the physical
method on real observations.}

\hypertarget{focused-acceptance-to-prepare-after-interface-review}{%
\subsection{Focused acceptance to prepare after interface
review}\label{focused-acceptance-to-prepare-after-interface-review}}

Before execution, announce the inputs, functions, actions and pass
criteria.

\begin{enumerate}
\def\labelenumi{\arabic{enumi}.}
\tightlist
\item
  Recover the independently calculated results in the synthetic example
  above.
\item
  Verify that BPM exclusions and nonfinite samples select exactly the
  same population for all means and both variances. Check ROI
  coordinates and counts.
\item
  Swap the members of either pair: gain and RON must be unchanged, while
  the corresponding difference median changes sign. Verify separately
  that subtracting a constant leaves variance and MAD unchanged.
\item
  Check constant images, too few accepted pixels, nonpositive gain
  denominators and invalid maps. Report unavailable values or errors as
  appropriate; preserve caller output on error and never modify the
  inputs.
\item
  Exercise the new call through the existing operation template using
  synthetic pairs whose members are not consecutive. Check failure
  handling and cleanup.
\item
  Verify the factor-of-two convention with a synthetic constant-gain
  relation. Curved examples must retain signal-dependent pair results
  rather than being silently presented as one constant gain.
  Sequence-level fitting remains later.
\item
  Run focused Mac C17 checks first and validate the changed path on
  Chapman afterward. Do not repeat unchanged qualification matrices or
  alter production installations or original data.
\end{enumerate}

\hypertarget{connection-to-the-full-linearity-design}{%
\subsection{Connection to the full linearity
design}\label{connection-to-the-full-linearity-design}}

The earlier
\href{detector-characterization-design-2026-11-02.pdf}{formula guide
(PDF)} and \href{detector-characterization-design-2026-11-02.tex}{LaTeX
source} remain the controlling design. This increment implements only
its paired gain/RON part.

The later shared shutter and linearity fit retains:

\[
Q_a(s)=s+c_{2,a}s^2+c_{3,a}s^3,
\qquad
Q_a(s_{rj})=L_r(t_j+\delta).
\]

\begin{longtable}[]{@{}
  >{\raggedright\arraybackslash}p{(\columnwidth - 4\tabcolsep) * \real{0.3333}}
  >{\raggedright\arraybackslash}p{(\columnwidth - 4\tabcolsep) * \real{0.3333}}
  >{\raggedright\arraybackslash}p{(\columnwidth - 4\tabcolsep) * \real{0.3333}}@{}}
\toprule\noalign{}
\begin{minipage}[b]{\linewidth}\raggedright
Symbol
\end{minipage} & \begin{minipage}[b]{\linewidth}\raggedright
Meaning
\end{minipage} & \begin{minipage}[b]{\linewidth}\raggedright
Unit
\end{minipage} \\
\midrule\noalign{}
\endhead
\bottomrule\noalign{}
\endlastfoot
\(a\) & Amplifier, or pixel in the later per-pixel mode & Identifier \\
\(r\) & Analysis region & Identifier \\
\(j\) & Exposure & Index \\
\(s\) & Bias-subtracted measured signal & ADU \\
\(s_{rj}\) & Signal for region \(r\) and exposure \(j\), subject to the
guide's pixel-model/aggregation rule & ADU \\
\(Q_a(s)\) & Linearized signal for model \(a\) & ADU \\
\(c_{2,a}\) & Quadratic coefficient of the inverse response &
\(\mathrm{ADU}^{-1}\) \\
\(c_{3,a}\) & Cubic coefficient of the inverse response &
\(\mathrm{ADU}^{-2}\) \\
\(L_r\) & Constant linearized illumination rate in region \(r\) &
ADU/s \\
\(t_j\) & Exposure duration established from the declared input-time
convention & s \\
\(\delta\) & One signed duration offset shared by the mosaic; positive
means longer exposure & s \\
\end{longtable}

The effective duration is \(t_j+\delta>0\). Keep the guide's
monotonicity, validated signal bounds, rank, conditioning, covariance
and aggregation requirements. No independent per-chip shutter offsets or
unrestricted lamp-drift parameters are introduced. The exact joint
numerical interface is a later review.

\hypertarget{selected-mac-omegacam-inputs}{%
\subsection{Selected Mac OmegaCAM
inputs}\label{selected-mac-omegacam-inputs}}

Root: \path{/Users/fselman/Pictures/OCAM/bpm}. Owner-provided listing;
files/pixels have NOT yet been independently inspected on the Mac. Ten
BIAS/r\_SDSS exposures precede twenty FLAT,DOME,DETGAIN/v\_STRM
exposures on 2017-03-13. Different bias/flat filters are recorded;
detector/readout compatibility still needs header confirmation.

Bias file timestamps
(\passthrough{\lstinline!OMEGA.2017-03-13T<timestamp>.fits!}):
09:50:05.624, 09:50:48.373, 09:51:31.403, 09:52:14.404, 09:52:57.394,
09:53:40.384, 09:54:23.384, 09:55:06.394, 09:55:49.384, 09:56:31.804.

\begin{longtable}[]{@{}
  >{\raggedright\arraybackslash}p{(\columnwidth - 8\tabcolsep) * \real{0.1875}}
  >{\raggedleft\arraybackslash}p{(\columnwidth - 8\tabcolsep) * \real{0.2500}}
  >{\raggedright\arraybackslash}p{(\columnwidth - 8\tabcolsep) * \real{0.1875}}
  >{\raggedright\arraybackslash}p{(\columnwidth - 8\tabcolsep) * \real{0.1875}}
  >{\raggedright\arraybackslash}p{(\columnwidth - 8\tabcolsep) * \real{0.1875}}@{}}
\toprule\noalign{}
\begin{minipage}[b]{\linewidth}\raggedright
Pair
\end{minipage} & \begin{minipage}[b]{\linewidth}\raggedleft
Time (s)
\end{minipage} & \begin{minipage}[b]{\linewidth}\raggedright
First flat timestamp
\end{minipage} & \begin{minipage}[b]{\linewidth}\raggedright
Second flat timestamp
\end{minipage} & \begin{minipage}[b]{\linewidth}\raggedright
Chronological flat indices
\end{minipage} \\
\midrule\noalign{}
\endhead
\bottomrule\noalign{}
\endlastfoot
P02 & 2 & 11:34:43.759 & 11:58:45.683 & 1,20 \\
P70 & 70 & 11:35:30.570 & 11:56:50.383 & 2,19 \\
P40 & 40 & 11:37:25.420 & 11:55:25.563 & 3,18 \\
P05 & 5 & 11:38:50.390 & 11:54:35.393 & 4,17 \\
P10 & 10 & 11:39:40.380 & 11:53:40.483 & 5,16 \\
P30 & 30 & 11:40:35.390 & 11:52:25.393 & 6,15 \\
P20 & 20 & 11:41:50.391 & 11:51:20.392 & 7,14 \\
P50 & 50 & 11:42:55.391 & 11:49:45.392 & 8,13 \\
P60 & 60 & 11:44:30.391 & 11:48:00.392 & 9,12 \\
P07 & 7 & 11:46:16.391 & 11:47:08.392 & 10,11 \\
\end{longtable}

This sequence has ten levels, not the design's representative 22-pair
annual sequence. Do not reject it solely on that count or promise all
products: admit each requested model using its actual rank, usable
signal span, saturation and illumination evidence. Preserve header
times; do not assume they are commanded times or already
shutter-corrected.

\end{document}
